% 04-walls.tex: схема как граница безопасности.
\section{Стены: схема как граница безопасности}
\label{sec:walls}

\subsection{Почему база данных не есть водопровод}

В большинстве систем база данных хранит то, что решило приложение. В Polaris решает база данных, а
приложение есть один клиент из нескольких: операторское веб-приложение, инструмент командной
строки, массовый загрузчик, консоль базы данных, восстановление из резервной копии. Правило,
живущее в коде приложения, связывает лишь тех клиентов, которые этот код выполняют. Правило,
закреплённое триггером, ограничением целостности или уникальным индексом, связывает каждого
клиента, переживает всякое восстановление и не может быть обойдено следующим вызывающим, потому
что у роли приложения нет привилегии определения данных, которой его можно было бы снять.
Управляющее предложение всей архитектуры звучит так: \emph{по построению}. Систему невозможно
эксплуатировать способом, нарушающим ограничение, потому что ограничение закрепляется до того, как
выполнится хоть строка кода приложения. \sImpl

Проектный принцип за этим старше Polaris: сделать недопустимые состояния непредставимыми, а не
просить разработчиков их не создавать. Событие проверки, записывающее идентификатор токена на
уровне раскрытия с нулевым разглашением, есть не ошибка приложения, которую надлежит отловить на
ревью. Это нарушение приватности, которое обязано быть нехранимым, и потому оно есть ограничение
целостности на строке.

\subsection{Десять ограничений и где живёт каждое}

Конституция называет десять жёстких ограничений в двух ярусах. Пять суть конституционные, гарантии
прав, которые нельзя ослабить без поправки, занесённой в журнал изменений; пять суть инженерные
инварианты, дисциплины, удерживающие первые пять истинными под нагрузкой и под атакой. Иерархия
глубиной в одну ступень, призвание, затем конституционные ограничения, затем инженерные
инварианты, и конституция гласит, что более глубокого механизма, ими управляющего, не существует.

\begin{table}[H]
\centering\footnotesize
\caption{Десять ограничений на \code{1.0.0-rc.7}, их ярус и объект, закрепляющий каждое. Таблица удерживается в согласии с кодом проверкой \code{check\_c1c10\_objects\_resolve}, которая роняет сборку, если названный объект перестал существовать.}
\label{tab:constraints}
\setlength{\tabcolsep}{4pt}\renewcommand{\arraystretch}{1.1}
\rowcolors{2}{tabstripe}{white}
\begin{tabularx}{\textwidth}{>{\bfseries}p{0.3in} >{\RaggedRight\arraybackslash}X >{\RaggedRight\arraybackslash}p{0.95in} >{\RaggedRight\arraybackslash}p{2.0in}}
\toprule
\rowcolor{tabheader}\thead{\#} & \thead{Ограничение} & \thead{Ярус} & \thead{Чем закреплено} \\
\midrule
C1 & След аудита работает только на добавление; история никогда не переписывается и не удаляется & Конституционное & Триггеры отвергают UPDATE и DELETE на каждой таблице аудита (\code{reject\_audit\_modification}) \\
C2 & Проверка с нулевым разглашением не хранит идентификатора токена & Конституционное & Двунаправленный CHECK \code{chk\_disclosure\_token\_consistency} \\
C3 & Человек держит не более одного ACTIVE-токена & Конституционное & Частичный уникальный индекс \code{uq\_one\_active\_per\_person} \\
C4 & Счёт неудачных входов атомарен & Инженерное & Однооператорный \code{UPDATE \ldots RETURNING} \\
C5 & Никаких встроенных сценариев; политика безопасности содержимого есть \code{script-src 'self'} & Инженерное & Заголовок ответа, проверяемый на каждом маршруте \\
C6 & Уровень раскрытия закрепляется на сервере; клиент не в силах повысить то, что узнаёт & Конституционное & Серверный код в паре с тестами сокрытия; стоящий за ним CHECK C2 \\
C7 & Никакой зашитой криптографии; алгоритмы суть строки в реестре & Инженерное & Внешний ключ на \entity{CryptographicAlgorithm} \\
C8 & Всякая карта и всякая сводка API ограничены & Инженерное & Жёсткие потолки в SQL-функциях и маршрутах \\
C9 & Утверждения об одновременности проверяются настоящими потоками & Инженерное & Многопоточные наборы тестов против живой базы данных \\
C10 & Личность не есть деньги; схема не несёт денежных полей & Конституционное & Структурное отсутствие, закреплённое проверкой \\
\bottomrule
\end{tabularx}
\end{table}

Репозиторий фиксирует также структурное утверждение о десятке: что набор замкнут, так что
одиннадцатое ограничение требует либо заменить одно из имеющихся, либо расширить рамку, и что
ограничения зависят друг от друга, так что удаление любого отзывается во всех остальных. Обход
зависимостей (уберите C10, и всякое иное ограничение защищает вместо этого систему финансовой
слежки; уберите C1, и закрепление всякого иного ограничения становится задним числом отрицаемым)
есть довод в пользу обращения с ними как с графом, а не как со списком. Афина, в
разделе~\ref{sec:cognition}, держит отображение каждого ограничения в точный живой механизм,
который его закрепляет, и проверка роняет сборку, если механизм уплывает. У трёх из десяти теперь
есть также формальная модель, проверяемая при каждой отправке: покрытие очистки для C1,
несвязываемость для C2 через две объединяющие данные доверяющие стороны и C3 при одновременных
выдаче и отзыве, и у каждой есть парная конфигурация, на которой она обязана упасть, так что ни
одна из трёх не пуста.

\figfit{0.635}{districts}{45 канонических таблиц \code{01\_schema.sql} на \code{1.0.0-rc.7} в виде
колец вокруг ядра версии 1. Бледная решётка в центре есть геометрия ядра с
рис.~\ref{fig:nucleus} в тех же пропорциях. Второе кольцо есть механика удостоверения: подписи,
разрешения, эпохи, восстановление, принуждение, ключ держателя, карта. Третье кольцо есть
полномочия, доверие и люди за стойкой: кто вправе подписывать по какому алгоритму, кто кому
доверяет, потолки отзыва, квоты, доверяющие стороны, регистры повторов и то, на чём держалась
регистрация. Четвёртое кольцо есть архивы, сторожевые башни и решения, сохранённые как записи:
журналы только на добавление, якоря, решение о сроках хранения в виде данных и события, которые
записывают всякое изменение органа, операторской учётной записи или доверяющей стороны. В
мигрированном развёртывании держится ещё семь таблиц (реестр миграций, операторские сессии и
учётные данные, журнал доступа к аудиту и три подготовленные таблицы Афины) и пять дочерних
разделов.}{fig:districts}

\subsection{Аудит как сама запись}

Четырнадцать таблиц следуют одному названному образцу, аудиту как самой записи: строка и её
история суть один и тот же предмет. Рядом с таблицей нет отдельного журнала событий, который можно
было бы править, пока таблица заперта. Изменение либо запрещено начисто, либо ограничено одним
полем, движущимся в одну сторону, вроде даты устаревания подписи или даты отзыва аттестации.
Каждая из четырнадцати закреплена триггером, поднимающим \code{insufficient\_privilege}, и
проверка роняет сборку, если хоть один триггер перестанет это делать. Сверх тех четырнадцати,
которые называет проектная запись, триггер добавления несёт на \code{1.0.0-rc.7} тридцать одна
таблица, и новые суть записи, которые теперь оставляют за собой решения об органах, операторах,
доверяющих сторонах и регистрациях; а начиная с \code{v9.412} C1 проверяется по одному разу на
каждую заявляющую его таблицу: управляемый таблицей набор тестов читает триггеры добавления из
системного каталога, а не из списка, который кто-то написал, так что новая таблица с добавлением
не может просто остаться непроверенной. \sImpl

\figside{town-map}{Обнесённый стеной город, с настоящей составной частью рядом с каждым городским
именем. Реестр в ратуше есть ядро удостоверения; стена есть схема; четверо ворот суть маршруты,
через которые утверждение входит и выходит; сторожевые башни суть журналы прозрачности, а
свидетели за стеной за ними наблюдают; дорога во второй город есть направленная аттестация,
привязанная к контексту; законы над городом связывают сам город. План описывает устройство.
Polaris федеративен: каждый орган есть свой собственный город.}{fig:town-map}

Следствие состоит в том, что ни один внешний ключ нигде в схеме не каскадирует удаление. Каскад
позволил бы событиям жизненного цикла удалённого токена исчезнуть вместе с ним. Всякая ссылка по
умолчанию есть \code{NO ACTION}, так что участник, на которого ссылается хоть одна строка аудита,
фактически неудаляем, а проверка просматривает каждый SQL-файл на обе формы каскада без единого
исключения. Единственная запись, намеренно не несущая никакого внешнего ключа, чтобы прекращение
договора с доверяющей стороной не могло стереть её историю, научила репозиторий кое-чему о его
собственной тестовой базе данных на \code{1.0.0-rc.7}: посев, перезагружающий образцовый мир,
перезапускал идентификаторы сторон с единицы и доходил до этой записи, не имея ни того ни другого,
так что первая сторона каждой перезагрузки наследовала события мёртвой стороны. Посев теперь
называет эту запись, а проверка вычисляет из собственных внешних ключей схемы, какие записи
перезагрузка оставила бы позади.

Роли приложения, \code{polaris\_app}, отозваны UPDATE и DELETE на таблицах с добавлением, и она не
держит привилегии определения данных. Единственный путь, удаляющий строки аудита, есть процедура
<<сначала в архив, потом очистка>>, которая сперва сверяет архив с его описью, отвергает всякую
отсечку внутри окна хранения и дописывает собственную контрольную строку в режиме добавления. Само
окно хранения есть данные с нижним пределом: ни одна строка политики не способна выразить срок
хранения короче года, потому что ограничение целостности этого не допускает, а понижение предела
есть изменение схемы, за которое отвечают на ревью. Призвание отвергает неограниченное хранение;
оно ровно так же отвергает хранение настолько короткое, что оно стирает запись. \sImpl

\subsection{Слой инвариантных проверок и учения, которые на него нападают}

Над схемой лежит плоский слой из 313 обыкновенных функций, каждая вида
\code{check\_*(repo\_root)}, которые читают репозиторий и возвращают вывод. Они запирают
непрерывную интеграцию: любой отказ роняет сборку. Каждая проверка сопряжена с тестом
обнаружения, который ломает образец ровно тем способом, какой проверка обещает поймать, и
утверждает, что проверка краснеет; а начиная с \code{v9.405} он утверждает ещё и то, что проверка
проходит на том же образце в целом виде, потому что тест, который только и делает, что
утверждает отказ, не установил, что он вообще что-нибудь обнаруживает. Проверка, не способная
обнаружить собственное нарушение, считается сломанной. Это и есть механизм, которым утверждения
прозы обращаются в запоры: величина, названная в README, маршрут, обязанный вернуть заголовок,
миграция, обязанная не каскадировать, учение, которое CI обязан прогнать. \sImpl

Версия 2 описала проверки и на этом остановилась. Между \code{v9.398} и \code{v9.437} репозиторий
задал этому слою тот вопрос, какой задаёт всему остальному, не \emph{проходит ли он}, а
\emph{заметил бы ли он}, и ответом часто было нет. Закомментирование строк, несущих искомые строки
проверки, что выглядит ровно как временно отключённая строка, оставило зелёными 71 проверку из
75, потому что вспомогательная функция, через которую читает каждая проверка, выдавала ей
комментарии вместе с кодом. Удаление каждого из 46 ограничений CHECK показало, что все до единого
несущие, и это была хорошая новость; удаление каждого из 37 триггеров нашло 14, которых никто не
заметил, и двенадцать из них были стражами добавления, которыми C1 \emph{и является}; удаление
каждого из 17 уникальных индексов нашло двенадцать непокрытых и остановилось на втором, потому что
в отсутствие \code{uq\_one\_active\_per\_person} набор тестов, существующий ради доказательства
C3, успел вставить и зафиксировать второй действующий токен для одного человека, и индекс уже
нельзя было перестроить. Удаление каждого из 59 отказов, которые делают шестнадцать хранимых
процедур, нашло 27 таких, что ни один тест по ним не затосковал бы, и сгущены они были ровно там,
куда триггер не видит: правило четырёх глаз, период выдержки, три канала и третий свидетель обряда
восстановления, предусловия необратимого стирания. Каждая из этих вещей теперь покрыта, и каждое
из этих измерений теперь есть учение, выполняемое при каждой отправке с отрицательным контролем,
так что результат <<ноль выживших>> отличим от стенда, который не прогнал ничего. На
\code{1.0.0-rc.7} таких учений двенадцать: над проверками, ограничениями, триггерами, процедурами,
свидетельствами с нулевым разглашением, контрактом соответствия, двумя эталонными верификаторами,
верификатором OpenID4VP, приложением, конституционными ограничениями, поверхностями, по которым
квантифицирует проверка, и перемещениями кода между модулями. Раздел~\ref{sec:senses} несёт эту
запись. \sImpl

\begin{layercard}{Карточка слоя: стены}
\qa{Что это}{PostgreSQL 16: 45 таблиц, 42 триггера, шестнадцать хранимых процедур, реализующих каждый сценарий со сменой состояния, таблицы аудита только на добавление, политика хранения в виде данных.}
\qa{Зачем оно существует}{Обещание, закреплённое лишь людьми, которые эксплуатируют систему, обещанием не является. Схема связывает каждого клиента, каждое восстановление и каждого инсайдера с подключением к базе данных, кроме владельца.}
\qa{Чему доверяет}{Движку базы данных и его владельцу; правильности ограничений в том виде, как они написаны, ради чего каждое теперь и проверяется мутацией.}
\qa{Чему отказывает}{Коду приложения, процессу ревью и дисциплине оператора как последней линии: приложение никогда не закрепляет конституционную гарантию само по себе. Тесту, которому ни разу не показали, как он падает.}
\qa{Что видит}{Всё, что система хранит. Поэтому слой приватности ограничивает то, что хранится, а не то, кому дозволено читать.}
\qa{Чего не видит}{Биометрические шаблоны, ключи на стороне держателя, регистрационные документы, содержимое обмениваемых сообщений или документов с отметкой времени: ничто из этого в базу данных не попадает.}
\qa{Если откажет}{Убранный триггер, пониженный предел хранения или каскад роняют сборку через проверки; ограничение, утраты которого набор тестов не заметил бы, роняет мутационное учение; у владельца базы данных с консолью возможностей больше, чем у любого триггера, о чём модель угроз говорит, а не умалчивает.}
\qa{Кем проверяется}{\code{check\_aor\_append\_only\_triggers}, \code{check\_aor\_privilege\_boundary}, \code{check\_no\_fk\_cascade}, \code{check\_retention\_engine}, \code{check\_append\_only\_tables\_are\_tested\_exhaustively}, \code{check\_constraint\_suite\_is\_mutation\_tested}, \code{check\_triggers\_are\_mutation\_tested}, \code{check\_unique\_rules\_are\_tested\_exhaustively}, \code{check\_procedure\_refusals\_are\_mutation\_tested}, \code{check\_seed\_restart\_resets\_dependent\_records}, самопроверки SQL.}
\qa{Сейчас или потом}{Сейчас.}
\qa{Внутри и снаружи}{Внутри: ядро удостоверения. Снаружи: криптографические свидетельства, потому что ограничение связывает лишь ту базу данных, которая его держит.}
\end{layercard}


\nextlayer{Ограничение связывает базу данных, а не копию утверждения, эту базу покинувшую. В тот
миг, когда удостоверение передано постороннему или историю органа читает аудитор, ничто из схемы с
ними не едет. Следующий слой заставляет утверждения везти доказательство с собой.}
